\documentclass[french]{article}
\usepackage{verbatim}
\usepackage{amsmath,amsfonts,amssymb}
\usepackage{pgfplots}
\usepackage[utf8]{inputenc}
\usepackage[french]{babel}
\usepackage{pgf,tikz,tkz-tab}
\usepackage{mathrsfs}
\usetikzlibrary{arrows}
\usepackage{pstricks}
%\pgfplotsset{compat=1.15}
\usepackage{geometry}
\geometry{hmargin=0.5cm,vmargin=0.5cm}


\begin{document}
\twocolumn
\begin{center}
\underline{\textbf{Multiple, divisiseur et nombre premier}}
\end{center}
\section*{Exercice I}
1)Donner les critères de divisibilité pour les nombres suivants: 
2; 3 et 7.

2)En sachant que $6=2\times 3$ donner un critère de divisibilité pour $6$.


3)À partir de la question précédente, déterminer les diviseurs des nombres suivant:
\begin{itemize}
\item[$\bullet$]42
\item[$\bullet$]1296 %42^2
\item[$\bullet$]1029 %7^3*3
\end{itemize}

\section*{Exercice II}
1)Compléter le tableau suivant:

\begin{tabular}[scale=1]{|l|c|c|c|} 
\hline
%\backslashbox{Type}{Variation sur $\mathbb{R}^+$}
Le nombre $\backslash$ appartient à & $\mathbb{N}$ & $\mathbb{Z}$ & traduction ensembliste \\
\hline
2 & Oui   & Oui & 2 $\in \mathbb{N}\cap\mathbb{Z}$  \\
\hline
-3 &  &  &  \\
\hline
5 &  &  &  \\
\hline
0 &  &  &  \\
\hline
1000 &  &  &  \\
\hline
-13 &  &  &  \\
\hline
\end{tabular}

2)Existe t-il un nombre qui appartient à $\mathbb{N}$ et qui n'appartient pas à $\mathbb{Z}$ ?

3)Existe t-il un nombre qui appartient à $\mathbb{Z}$ et qui n'appartient pas à $\mathbb{N}$ ?

\section*{Exercice III}
Soit un carré de coté $c \in \mathbb{N}$ et d'aire $A \in \mathbb{N}$.\\

1) Rappeler quelle est l'aire d'un carré de coté $c$.

2) Calculer l'aire $A$ pour les valeurs de $c$ suivantes: $1; 2; 3; 4; 5; 6; 7; 8; 9; 10$

3) Que remarque t-on sur l'aire $A$ si $c$ est paire?

4) Que remarque t-on sur l'aire $A$ si $c$ est impaire?

2) Est-il possible que $A=7$?

3) Est-il possible que $A=5$?

4) En sachant que $A\leq 10$, lister toutes les valeurs possibles de $c$.

\section*{Exercice IV}

Soit un rectangle de largeur $l=2$ et de largeur $L \in \mathbb{N}$ et d'aire $A$.

1) Rappeler quelle est l'aire d'un rectangle de largeur $l$ et de longueur $L$.

2) Est-il possible que l'aire $A$ soit impaire ?

3) Est-il possible que $A=27$ ? Si oui, déterminer $L$.

4) Est-il possible que $A=86$ ? Si oui, déterminer $L$.
 

\section*{Exercice V}
Répondre par vrai ou faux:

-Il n'existe pas de nombre pair multiple de 3.

-Si un nombre est un multiple de 3 et de 4 alors il est multiple de 7.

-Soit $n\in \mathbb{N}$, est-ce que $3n+9$ est un multiple de 3?  

-0 est un nombre pair.

-0 est premier.

-0 est divisible par 3.


\section*{Exercice VI}
Recopier les phrases suivantes en complétant par "multiple" ou "diviseur":

250 est un ... de 50

1000 est un ... de 1

30 est un ... de 30

5 est un ... de 55

3 est un ... de 333 666 666

250 est un ... de 50

\section*{Exercice VII}
Combien il y a t'il de multiple de 13 entre 1 et 1000?

Combien il y a t'il de multiple de 13 entre 1 et 500?

Combien il y a t'il de multiple de 13 entre 500 et 1000? \vspace{1.8cm}



\section*{Exercice VIII}
Soit ici $\phi$ une fonction tel que $D_f=\mathbb{P}$, l'ensemble des nombres premiers que l'on range par ordre croissant, on associe à ce nombre premier la position de celui-çi dans l'ensemble $\mathbb{P}$.


\underline{Exemple:} 

$\phi(2)=1$ car $2$ est le \underline{premier} nombre premier,

 
$\phi(3)=2$ car $3$ et le \underline{deuxième} nombre premier etc.

1) Trouver $\phi(11)$ 

2) Trouver $p \in \mathbb{P}$ tel que $\phi(p)=10$

\end{document}


